\documentclass[../main.tex]{subfiles}
% Do not include packages here
\begin{document}

\chapter{Estimation of Parameters}

The general problem of parameter estimation may be sketched as follows:

From a restricted number of observations,
assumed to constitute a random sample,
one wants to gain some knowledge about the underlying population or universe
from which the sample emanated.
The mathematical form of the parent distribution may be well-defined,
but it involves a certain number of parameters
whose numerical values are not known.
The measurements should therefore be used to extract the largest possible amount of information about the parameters,
specifically,
the experimental sample should supply numbers which could be said to represent the numerical values of the parameters.

We have already in Chapter 7 seen examples on interval estimation of the parameters in the normal distribution
In the forthcoming chapters,
we shall study different methods for estimating unknown parameters in general probability distributions
These methods
(the Maximum-Likelihood-, Least-Squares- and moments methods)
all produce point estimates,
that is, some definite numbers for the parameters.
The methods also give measures for the uncertainties in the estimates
and thereby express the confidence that can be attached to the numerical results.

In this chapter,
we shall take up various general aspects of parameter estimation
by discussing in some detail a few of the criteria that should be fulfilled by good and acceptable estimators.
Although these criteria will be applied to the specific point estimation methods described in Chapters \ref{ch:9}-\ref{ch:11},
the discussion in the following sections will mainly be of a formal nature.
The student who requires just a superficial knowledge of these rather theoretical features
may therefore be satisfied with reading only the first two sections of this chapter.

\section{Definitions}

The term estimator denotes in the following a function of the observations,
or the method or prescription used to find a value for an unknown parameter.
By an estimate we mean the numerical value of the parameter obtained with the estimator for a particular set of observations.
If the parameter is $\theta$, its estimate is denoted by $\hat{\theta}$.
The term statistic was introduced in Sect. \ref{sec:7.1}
as a function of one or more random variables that does not depend on any unknown parameters.
In the general case,
we will let the statistic $t = t(x_1,x_2,\ldots,x_n)$
be an estimator of the unknown $\theta$ or of some function of $\theta$.
If a continuous or discrete population has the probability distribution $f(x; \theta)$,
the \emph{likelihood} of the observations $x1, x2, \ldots, x_n$ for a specific $\theta$ is given by
%
\begin{equation}
    L(x_1, x_2, \ldots, x_n | \theta) = \prod_{i=1}^n f(x_i ; \theta)
    \label{eq:8.1}
\end{equation}
%
The product expresses the joint conditional probability for obtaining the
measurements $x_1, x_2, \ldots, x_n$ given $\theta$.
We shall later also think of $L(x_1, x_2, \ldots, x_n | \theta)$ as a function of $\theta$,
calling it a \emph{likelihood function}.
We will use the symbol $L$ or $L(\vec{x} | \theta)$ to denote the likelihood
(i.e, the number expressing the joint probability)
as well as the likelihood function
(i.e. the function of $\theta$).
The third interpretation regards $L$ as a joint probability density function of observable quantities $x_1, x_2, \ldots, x_n$, for a given $\theta$.
Hopefully, in the following,
the precise meaning of the symbol $L$, or $L(\vec{x} | \theta)$,
will be clear from the context in which it appears.

\section{Properties of estimators}

Observations are random variables.
Any function of the observations will also be a random variable,
which may take on a variety of values.
If a statistic $t = t(x_1,x_2,\ldots,x_n)$ is used as an estimator for the parameter $\theta$,
it will therefore give rise to a distribution of estimates $\hat{\theta}$.
The individual estimates obtained are of less interest than their overall distribution,
because this distribution will reflect the quality of the estimator when it is used many tines.
We will therefore judge the merits of an estimator from the characteristics of the distribution of its estimates.

A good estimator should in the long run produce estimates
which do not systematically deviate from the true parameter value,
and its accuracy should increase with the number of observations.
Frequently, there are several estimators which fulfill these requirements
and hence can reasonably be thought of for estimating an unknown parameter.
If so, one estimator can be said to be superior to the others
if its distribution of estimates shows the best "concentration" about the true parameter value. 
"Concentration" may for this purpose be expressed by giving the variance as a measure of the spread of the distribution about its central value.

In the forthcoming sections,
we will discuss the following optimum properties that are desired for good estimators:
consistency, unbiasedness, minimum variance, efficiency and sufficiency.
Only rarely will the conceivable estimators for a parameter possess all the good properties.
One may therefore have to choose between them,
and in each specific case decide which of the ideal properties that can be abandoned,
taking also practical considerations into account.

\section{Consistency}

It is intuitively clear that a desirable property of an estimator
is that its estimates converge towards the true parameter value
when the number of observations is increased.
Such a property is called consistency.

Mathematically, the estimator $t$ can be said to be consistent if it converges in probability to $\theta$.
If the estimate $\hat{\theta}_n$ is obtained from a sample of size $n$,
then, given any positive $\epsilon$ and $\delta$,
an $N$ should exist such that
%
\begin{equation}
    P(|\hat{\theta}_n - \theta| > \epsilon) < \delta
    \label{eq:8.2}
\end{equation}
%
for all $n > N$.
In words,
the estimator is consistent if,
given any small quantity $\epsilon$,
we can find a sample size $N$ such that,
for all larger samples,
the probability that $\hat{\theta}_n$ differs from the true value by more than $\epsilon$
is arbitrarily close to zero.

As an example,
we know from the Law of Large Numbers (Sect. \ref{3.10.4}) that
the arithmetic mean of a sample of $n$ measurements
from a population with mean $\mu$ and finite variance
will converge towards $\mu$ as $n$ becomes large,
%
\begin{equation}
    t
    = \bar{x}
    = \frac{1}{n} \sum_{i=1}^n x_i
    \xrightarrow[n \to \infty]{} \mu
\end{equation}
%
The sample mean is therefore a consistent estimator of the population mean.

\textbf{Exercise 8.1:}
Show explicitly that
the mean $\bar{x}$ of a sample of size $n$
from the normal population $N(\mu,\sigma^2)$
converges in probability to $\mu$.

\section{Unbiasedness}

Consistency is a property
that describes the behavior of an estimator when the sample size increases to infinity,
but says nothing about its behavior for finite data sets.
For example,
in the last section,
we just saw that the sample mean $\bar{x}$ is a consistent estimator of the population mean $\mu$,
but so would also be
%
\begin{equation}
    t' = \bar{x'} = \frac{1}{n-a} \sum_{i=1}^n x_i
\end{equation}
%
where $a$ is any fixed number.
Why do we prefer one to the other?

Unbiasedness is an estimator property defined for finite sets of observations.
This property is possessed by estimators whose estimates are not systematically shifted from the true parameter value
but centered around this value for all sample sizes.
As a measure of the centre of a distribution,
one can use the conventional mean value.
Mathematically, the property unbiasedness is therefore defined if,
for all sample sizes $n$,
the expectation value of the estimator $t$ is equal to the true parameter value $\theta$,
%
\begin{equation}
    E(t) = \theta
    \label{eq:8.3}
\end{equation}
%
if
%
\begin{equation}
    E(t) = \theta + b
    \label{eq:8.4}
\end{equation}
%
and $b$ is different from zero, the estimator is biased.
The bias term $b = b(\theta)$
will for all reasonable estimators be of order $\frac{1}{n}$ or smaller compared to $\theta$.

It is rather trivial to see that the sample mean $X$ is an unbiased
estimator of the population mean $\mu$ whenever the latter exists.
It is also seen that the estimator $\bar{x'}$ above
will be a biased estimator of $\theta$ for all $a$ different from zero.

Consistency and unbiasedness are independent estimator qualities,
as neither property implies the other.
It is generally accepted that consistency is more important than unbiasedness,
partly because bias can often be corrected for.
A consistent estimator whose asymptotic distribution has a finite mean will always be asymptotically unbiased.

\subsection{Example: $s^2$ as an estimator of $\sigma^2$}
\subsection{Example: Estimator of the third central moment}

\textbf{Exercise 8.2:}
In the binomial distribution $B(k | n,p) = \binom{n}{k} p^k (1-p)^{n-k}$,
where $k=0, 1, \ldots, n$,
show that the unbiased estimator of $p$ and $p^2$ are, respectively,
$k/n$ and $k(k-1)/n(n-1)$.

\section{Minimum variance and efficiency}

The requirements of consistency and unbiasedness do not uniquely determine how to choose a good estimator.
One finds, for instance, that both the sample mean and the sample median
are consistent and unbiased estimators of the location of a normal population with known variance.
However, as it can be shown that the variance of the mean is smaller than the variance of the median,
the mean is regarded as a better estimator of the central value.
It seems natural, therefore, to use the spread in the estimates as a measure for the acceptability of an estimator.
For most distributions encountered in practice,
the second central moment, or the variance, will be a good measure of the concentration of the estimates;
this is especially so for the many cases where this distribution is approximately normal.

Under fairly general conditions,
there exists a lower bound on the variance of the estimates derived from an estimator.
This lower bound is easily established when considering the likelihood function defined by \eqref{eq:8.1}.
We shall assume that
the first two derivatives of $L(\vec{x}|\theta)$ with respect to $\theta$ exist for all $\theta$,
and that the range of $x$ is independent of $\theta$.
Given an estimator $t$ of some function of $\theta$, say $\tau(\theta)$,
we define its bias $b(\theta)$ by the relation (compare \eqref{eq:8.4}),
%
\begin{equation}
    E(t)
    = \int \ldots \int
        t(\vec{x}) \,
        L(\vec{x} | \theta) \,
        d\vec{x}
    = \tau(\theta) + b(\theta)
    \label{eq:8.5}
\end{equation}
%
With the assumption that the range of $x$ is independent of $\theta$,
the differentiation of \eqref{eq:8.5} with respect to $\theta$ gives
%
\begin{equation}
    \int \ldots \int t \, \pderiv{L}{\theta} \, d\vec{x}
    = \int \ldots \int t \, \pderiv{\ln L}{\theta} \, L \, d\vec{x}
    = \pderiv{\tau}{\theta} + \pderiv{b}{\theta}
    \label{eq:8.6}
\end{equation}
%
Since $L$ is the joint probability of the observations,
%
\begin{equation}
    \int \ldots \int L d\vec{x} = 1
    \label{eq:8.7}
\end{equation}
%
we find by differentiating this relation with respect to $\theta$
%
\begin{equation}
    \int \ldots \int \pderiv{L}{\theta} \, d\vec{x}
    = \int \ldots \int \pderiv{\ln L}{\theta} \, L \, d\vec{x} 
    = 0
    \label{eq:8.8}
\end{equation}
%
Multiplying \eqref{eq:8.8} by $\tau(\theta)$
and subtracting the result from \eqref{eq:8.6} leads to
%
\begin{equation}
    \int \ldots \int [t - \tau(\theta)] \, \pderiv{\ln L}{\theta} \, L \, d\vec{x}
    = \pderiv{\tau}{\theta} + \pderiv{b}{\theta}
    \label{eq:8.9}
\end{equation}
%
By applying the Schwarz inequality to the integral,
we obtain the formula
%
\begin{equation}
    \left( \int \ldots \int \big[t - \tau(\theta)\big]^2 \, L \, d\vec{x} \right)
    \left( \int \ldots \int \pderiv{\ln L}{\theta}^2 d\vec{x} \right)
    \geq \left( \pderiv{\tau}{\theta} + \pderiv{b}{\theta} \right)^2
    \label{eq:8.10}
\end{equation}
%
which can be written as
%
\begin{equation}
    V(t)
    \geq \frac
    {\left( \pderiv{\tau}{\theta} + \pderiv{b}{\theta} \right)^2}
    {E\left[\left(\pderiv{\ln L}{\theta}\right)^2\right]}
    \label{eq:8.11}
\end{equation}
%
This fundamental inequality for the variance of an estimator is often referred to as the Cramer-Rao inequality.
When $L$ satisfies the aforementioned regularity conditions,
one can prove the following relation,
%
\begin{equation}
    E\left[\left(\pderiv{\ln L}{\theta}\right)^2\right]
    = E\left[-\frac{\partial^2 \ln L}{\partial \theta^2}\right]
    \label{eq:8.12}
\end{equation}
%
Thus an alternative form of the Cramer-Rao inequality,
which is sometimes easier to evaluate, is
%
\begin{equation}
    V(t)
    \geq \frac
    {\left( \pderiv{\tau}{\theta} + \pderiv{b}{\theta} \right)^2}
    {E\left[-\frac{\partial^2 \ln L}{\partial \theta^2}\right]}
    \label{eq:8.13}
\end{equation}
%
The lower limit of the variance implied by \eqref{eq:8.11}, \eqref{eq:8.13} is called the minimum variance bound, MVB,
and an estimator attaining this limit is called an MVB estimator,
or more often, and efficient estimator.

From the derivation above,
it is realized that the variance of an estimator will attain the MVB
if the Schwarz inequality applied to \eqref{eq:8.9} becomes an equality.
The necessary and sufficient condition for this
is that $[t-\tau(\theta)]$ is linearly related to $\pderiv{\ln L}{\theta}$ for all sets of observations;
we write
%
\begin{equation}
    \pderiv{\ln L}{\theta} = A(\theta) \left[ t - \tau(\theta) - b(\theta) \right]
    \label{eq:8.14}
\end{equation}
%
From \eqref{eq:8.13}, one then gets a simple formula for the MVB,
%
\begin{equation}
    V(t) = \frac
    {\left( \pderiv{\tau}{\theta} + \pderiv{b}{\theta} \right)^2}
    {A(\theta)}
    \label{eq:8.15}
\end{equation}
%
Efficient estimators exist only for the limited class of problems for which the likelihood function satisfies the condition \eqref{eq:8.14}.
Most estimators will have a variance larger than the MVB.
We therefore define the efficiency of an estimator as the ratio between the MVB and the actual variance $V(t)$ of the estimator,
%
\begin{equation}
    \text{Efficiency(t)} = \frac{\text{MVB}}{V(t)}
    \label{eq:8.16}
\end{equation}
%
It is interesting to note that
for distributions where the regularity conditions do not hold
the minimum attainable variance may be smaller or larger than the MVB.

For the particular case where $t$ is an estimator of the parameter $\theta$ itself,
we have $\pderiv{\tau}{\theta} = 1$,
and the Cramer-Rao inequalities \eqref{eq:8.11} and \eqref{eq:8.13} become, respectively,
%
\begin{align}
    V(t)
    &\geq \frac
    {\left( 1 + \pderiv{b}{\theta} \right)^2}
    {E\left[\left(\pderiv{\ln L}{\theta}\right)^2\right]}
    \label{eq:8.17}
    \\
    V(t)
    &\geq \frac
    {\left( 1 + \pderiv{b}{\theta} \right)^2}
    {E\left[-\frac{\partial^2 \ln L}{\partial \theta^2}\right]}
    \label{eq:8.18}
\end{align}
%
The efficiency condition \eqref{eq:8.14} then reads
%
\begin{equation}
    \pderiv{\ln L}{\theta} = A(\theta) \left[ t - \theta - b(\theta) \right]
    \label{eq:8.19}
\end{equation}
%
simplifying the MVB formula \eqref{eq:8.15} to
%
\begin{equation}
    V(t) = \frac
    {\left( 1 + \pderiv{b}{\theta} \right)^2}
    {A(\theta)}
    \label{eq:8.20}
\end{equation}
%

\textbf{Exercise 8.3:}
Prove \eqref{eq:8.12}.
(Hint: Differentiate \eqref{eq:8.8} with respect to $\theta$).

\subsection{Example: Estimator of the mean in the Poisson distribution}
\subsection{Example: Estimators of the mean in the normal p.d.f.}
\subsection{Example: Estimators of $\sigma^2$ and $\sigma$ in the normal p.d.f.}

\section{Sufficiency}

\subsection{One-parameter case}

An estimator $t$ is said to be \emph{sufficient}
if it exhausts all information in the observations $x_1, x_2, \ldots, x_n$ regarding the parameter $\theta$.
For the normal distribution,
we have used the sample mean $\bar{x} = \frac{1}{n} \sum_i x_i$
as an estimator of the population mean $\mu$.
No extra knowledge on $\mu$ can be gained from other functions of the observations,
such as $\sum_i |x_i|$, $\sum_i x_i^2$, etc.;
the statistic $\bar{x}$ is then a sufficient estimator for $\mu$.
Actually, any function of $\bar{x}$ provides a sufficient statistic for $\mu$ in the normal distribution.
Therefore,
to choose between the different functions,
one may have to test also the qualities of consistency, unbiasedness, and efficiency.

To be more precise,
let us consider the likelihood function when the p.d.f. is $f(x;\theta)$.
Suppose that $L$ can be factorized to give
%
\begin{equation}
    L(\vec{x}|\theta)
    = \prod_{i=1}^n f(x_i;\theta)
    = G(t|\theta) \, H(\vec{x})
    \label{eq:8.25}
\end{equation}
%
where the function $G$ involves the statistic $t$ and the parameter $\theta$,
and $H$ is independent of $\theta$,
being a function of the observations $\vec{x}$ only.
Since $G$ only has reference to the data via $t = t(x_1,x_2,\ldots,x_n)$,
and is a known number for the given sample,
$t$ must supply all the available information in the data regarding $\theta$.
It follows that whenever the likelihood function can be written in the form of \eqref{eq:8.25},
$t$ is a sufficient statistic for the parameter $\theta$.

One can show that a necessary condition for a sufficient statistic to exist is that
the p.d.f. belongs to the exponential family,
defined by
%
\begin{equation}
    f(x|\theta)
    = \exp \left[ B(\theta)C(x) + D(\theta) + E(x) \right]
    \label{eq:8.26}
\end{equation}
%
where $B,C,D,E$ are functions of the indicated arguments.

For the restricted class of p.d.f's satisfying \eqref{8.26},
one sees that the factorization requirement of \eqref{8.25}
implies that a sufficient statistic must be expressed by the function $C(x)$,
%
\begin{equation}
    t = \sum_{i=1}^n C(x_i)
    \label{eq:8.27}
\end{equation}
%
It is easily shown that efficient estimators are always sufficient.
To see this,
we take logarithms on both sides of \eqref{8.25}
and differentiate with respect to $\theta$,
getting
%
\begin{equation}
    \pderiv{\ln L}{\theta} = \pderiv{\ln G}{\theta}
    \label{eq:8.28}
\end{equation}
%
We see that the efficiency condition \eqref{8.14}
is just a special case of the more general \eqref{8.28}, with
%
\begin{equation}
    \pderiv{\ln G}{\theta} = A(\theta) \left[ t - \tau(\theta) - b(\theta) \right]
    \label{eq:8.29}
\end{equation}
%
Under the regularity conditions specified for the likelihood function in Section \ref{sec:8.5},
there is among all sufficient estimators for $\theta$ only one efficient estimator.
If the p.d.f. belongs to the exponential family
and the range of $x$ is independent of $\theta$,
it can be shown that sufficient statistics always exist for $\theta$.
However, there will be just one sufficient statistic which will satisfy \eqref{8.29}
and thus estimate some function $\tau(\theta)$ with variance equal to the MVB;
compare the example of Section \ref{sec:8.5.3}.
Furthermore, for large samples, any function of a sufficient statistic will be an MVB estimator.

\textbf{Exercise 8.8:}
Verify that the Poisson distribution
$P(x;\theta) = \frac{1}{x!} \theta^x e^-\theta$
belongs to the exponential family.
Compare Section \ref{sec:8.5.1}.

\textbf{Exercise 8.9:}
Show that the Cauchy p.d.f
$f(x;\theta) = \frac{1}{\pi} \frac{1}{1 + (x-\theta)^2}$
does not have a sufficient estimator of $\theta$.
Compare Exercise \ref{ex:8.4}.

\subsection{Example: Single sufficient statistics for the normal p.d.f.}
\subsection{Extension to several parameters}
\subsection{Example: Jointly sufficient estimators for $\mu$ and $\sigma$ in N($\mu$, $\sigma^2$)}

\end{document}
